% ============================================
% DOC-27 — A Origem (1853) (edição revista R1075)
% ============================================
\fragtitle{doc}{D}{DOC}{27}{A Origem (1853)}
\markboth{DOC-27 --- A Origem (1853)}{}

\noindent \textit{Arquivo Público do Ceará --- documento avulso --- século XIX. Anexado por Dra. Lúcia Mendes em 12 de fevereiro de 2026, com a anotação: ``Isto prova que o diário não é o começo. O diário é só o eco.''}

\noindent \textbf{Documento:} carta do Padre José Inácio de Sousa ao Presidente da Província do Ceará, Comissão de Socorros Públicos. \textbf{Data:} 14 de março de 1853. \textbf{Origem:} Freguesia de Nossa Senhora dos Remédios, Rio Quixeramobim. \textbf{Estado:} papel manchado por umidade; a borda inferior apresenta uma substância escura, não identificada, que atravessou todas as folhas. Análise química inconclusiva.

\vspace{0.5em}
\noindent Excelentíssimo Senhor Presidente da Província,

\noindent Tomo a liberdade de escrever-lhe porque os vivos já não me ouvem e os mortos são muitos para enterrar sozinho.

\noindent Estou aqui desde novembro do ano passado, destacado para administrar os sacramentos aos retirantes que chegam à margem do rio. Recebi duzentas almas em janeiro. Duzentas e quarenta em fevereiro. Neste mês de março, já perdi a conta. Chegam aos magotes, famílias inteiras reduzidas a pó ambulante. Os olhos deles são a primeira coisa que morre. Antes do corpo. Ficam abertos, mas não veem. Como se a seca tivesse entrado por eles e secado a alma por dentro.

\noindent Os que morrem são levados para uma vala a meia légua da freguesia. Não tenho cemitério consagrado. Não tenho madeira para caixões. Não tenho pano para mortalhas. Tenho apenas uma cova comprida, aberta na terra seca, onde os corpos são deitados em camadas. Terra sobre eles. Mais corpos. Mais terra. Os urubus aprenderam o caminho. Fazem um círculo no céu que não se desfaz. Nem de noite. Ficam empoleirados nas árvores ao redor, esperando.

\noindent E é sobre isto que preciso escrever-lhe.

\noindent Há três noites, não consigo rezar o ofício das almas sem sentir uma presença. Não uma presença espiritual, como as que estou acostumado a discernir em trinta anos de sacerdócio. É uma presença física. Ela chega com o cheiro.

\noindent O cheiro vem da vala.

\noindent O cheiro comum da morte eu conheço bem. Este é doce. Mais doce que incenso, mais doce que flor de mandacaru. Um cheiro que entra pelas narinas e se instala na garganta como um líquido grosso. Os retirantes sentem também. Eles se viram na direção da vala e ficam olhando. Não choram. Não rezam. Apenas olham. Como se esperassem que algo fosse sair de lá.

\noindent Eu mesmo fui até a vala na segunda noite. Levei uma lanterna, água benta, um crucifixo. Coisas de padre. Coisas que sempre funcionaram. Cheguei à borda e iluminei o fundo.

\noindent O que vi não eram corpos.

\noindent Eram restos. Braços e pernas saindo da terra como raízes. Mãos abertas para cima, como se pedissem. Uma mão de criança, pequena, com os dedos retorcidos, apontava para o céu noturno como se acusasse alguém.

\noindent O cheiro era mais forte ali.

\noindent Ajoelhei para benzer a vala. Comecei o \textit{Requiem Aeternam}. Na terceira palavra, senti que a terra vibrou. Não um tremor. Uma pulsação. Como se o chão tivesse um coração.

\noindent Pressionei a mão direita contra a terra. Ela estava morna. Não quente do sol. Morna de dentro. A temperatura do corpo de um doente. A terra estava febril.

\noindent Enterrei a mão até o punho. A terra era macia, úmida --- coisa impossível neste sertão seco. Senti algo roçar meus dedos. Algo liso. Mole. Algo que se moveu ao toque, como se estivesse vivo.

\noindent Tirei a mão. A terra fechou sozinha sobre o buraco.

\noindent Um homem que estava perto --- um retirante chamado Vicente, que perdeu a mulher e três filhos na semana passada --- olhou para mim e disse:

\noindent ``Seu padre, a vala tá com fome.''

\noindent Excelência, sei o que Vossa Senhoria pensará ao ler isto. Que o isolamento e a fome abalaram meu juízo. Que o calor e a mortandade turvaram minha razão. Não sei se é verdade. Talvez seja. Mas preciso registrar o que vi, para que, se eu morrer aqui --- e acho que vou morrer aqui ---, alguém saiba que esta vala não é uma vala comum.

\noindent No terceiro dia, acordei com a palma da mão direita marcada. Uma marca escura, arredondada, no centro da palma. Parece uma queimadura, mas não dói. A pele está íntegra. A marca apenas está lá. Como se alguém tivesse pressionado um selo quente contra minha carne enquanto eu dormia.

\noindent Tentei lavar. Não sai. Tentei rezar. A marca não some.

\noindent Talvez seja só a terra que entrou na minha pele. Talvez seja a úlcera que está começando. Talvez seja o juízo que se vai.

\noindent Mas, esta noite, quando fechei os olhos, ouvi uma voz. Não uma voz que vem de fora. Uma voz que vem de dentro da terra. Ela repetia a mesma palavra, sem parar, como quem aprende a falar:

\noindent ``Fome.''

\noindent ``Fome.''

\noindent ``Fome.''

\noindent A voz de uma criança. Mas nenhuma criança sobreviveu para falar.

\noindent Peço licença para encerrar esta carta. A lanterna está baixando.

\noindent Que Deus tenha misericórdia de nós.

\vspace{0.5em}
\noindent Padre José Inácio de Sousa --- Freguesia de Nossa Senhora dos Remédios --- 14 de março de 1853.

\vspace{0.8em}
\noindent \textit{Nota do Arquivo Público do Ceará: a carta foi encontrada em 1916, dentro de um livro de tombo da freguesia, durante a reforma do prédio. O sacerdote faleceu em 15 de março de 1853, um dia após redigir esta carta. A causa da morte foi registrada como ``fome voluntária'': o padre teria deixado de se alimentar, cedendo sua ração aos retirantes. O corpo foi enterrado na mesma vala que ele descreve.}

\vspace{0.5em}
\noindent \textit{Anotação a lápis, na margem da ficha de catalogação de 1916, em caligrafia não identificada:} ``O cheiro doce nunca saiu daqui. Transferir para o Arquivo Público Mineiro. Processo 1.260.''

\vspace{1em}
\noindent \textbf{Julho de 2026 --- Belo Horizonte. Último registro profissional de Dra. Lúcia Mendes.}

\noindent \textbf{Documento interno --- autoencaminhamento para afastamento clínico.} Profissional: Lúcia Mendes, CRM-MG 28.391. Data: 28 de julho de 2026.

\noindent Em janeiro de 2026, tive contato com o diário atribuído ao paciente 1.260, Joaquim Antônio Correia, no contexto de pesquisa em psicologia forense histórica. O material permaneceu sob guarda intermitente, com falhas de lacre e registro.

\noindent Nas semanas seguintes, observei: alteração subjetiva de coloração no olho direito, sem exame oftalmológico conclusivo; percepção recorrente de odor adocicado sem fonte ambiental; sonhos repetitivos com vala, enfermaria e formulário em branco; hipervigilância noturna; impulso de retornar ao diário mesmo após decisão formal de encerramento.

\noindent Hipóteses: reação aguda ao material traumático; sugestão intensificada pela exposição documental; identificação contratransferencial com o paciente; fenômeno dissociativo leve, a confirmar por avaliação externa.

\noindent Medidas: suspensão de entrevistas, escrita acadêmica e circulação de cópias; encaminhamento do original para custódia judicial; solicitação de avaliação por profissional sem contato prévio com o caso; recomendação de que o diário não seja lido individualmente.

\noindent \textit{Anotação não protocolada: a cópia impressa deste documento apresenta, ao final, a numeração ``1.262'' escrita a lápis ao lado do nome da profissional. O traço não corresponde à assinatura principal. Abaixo dele há uma linha vazia, sem identificação.}

\noindent Lúcia Mendes --- CRM-MG 28.391. Registro encerrado às 3h14.

> **↪ Rotas:** **DOC-01** · **LUC-02** · **EPI-01**
